% !TeX root = ../papers/nc-p3.tex
\section{非全射・非被約を含む一辺の二元射影の全面排除}
\label{sec:deficient-binary-quintic-pencil}
\begin{theorem}\label{thm:deficient-binary-quintic-pencil}
標準Koszul左右AS正則な四窓では、一文字を直線で割ったとき、射影された一辺の関係空間が
$P_5^*\subset P_2^*\otimes P_3^*$ の乗法双対像に含まれる基底は存在しない。
像の次元に条件はなく、辺の向きはどちらでもよい。点グラフの次元・被約性も任意である。
\end{theorem}
\begin{proof}[捩れ評価における二つの端支持]
像が六次元の場合は定理\ref{thm:binary-edge-basepoints}による。
真の部分空間 $H$ の場合は定理\ref{thm:deficient-binary-support}により
$s=\dim H$ は4か5、隣接評価の捩れ長 $t$ は2か3で $W=M_5$ である。
任意の捩れの台の点 $x$ における余核の剰余体商から、第四文字 $l_x\in W$ を得る。
その三テンソルは圧縮後 $\operatorname{ev}_x^{(5)}\otimes v_x$ に等しく、
圧縮前の左端支持は $D_x=\langle z,\operatorname{ev}_x^{(2)}\rangle$ に入る。
正則窓での端支持下界から丁度 $D_x$ である。同一点で二つの独立な剰余体商があれば
同じ二次元端支持を与えるので、定理\ref{thm:all-endpoint-ranks-two}に反する。
従って余核は次数 $t$ の有効因子 $T$ 上可逆である。

$s=5$ とすると、指定の辺の射影核 $z\otimes S$ は一次元である。
$w(l_x)$ は圧縮後の評価の一つの持上げと $v_x$ の積に、
この射影核と一つの末尾ベクトルの積を加えたものである。よって右端支持も高々二次元となる。
両端とも丁度二次元なら $l_x$ が前点・後続点の双方の直線を持ち、
定理\ref{thm:no-two-sided-centers}に反する。これで階数5が除かれる。
\end{proof}
\begin{proof}[二次式による両端の同時切断]
$s=4$ とする。$T$ 上で束を自明化し、$T^*$ を $H^0(\mathcal O_T)^*$ の
各次数の評価逆系と書く。前節の二つの束列から $T^*$ が $M_n$ に入り、
その上の多項式による縮約は $T$ の環の乗法の双対である。
$\dim(W/T^*)=4-t$ である。また $H\subset P_5^*$ は通常のスカラーの評価空間
$H^0(\mathcal O_T(5))^*$ を含む。実際 $V$ の文字で $T$ 上いたるところ単元となる
一つを選べる。各閉点で剰余値が非零の条件は真の超平面の補であり、文字は層を生成するからである。
この文字と $W$ 内の $T^*$ の全てとで圧縮窓を評価すると、$T$ 上の全五次汎関数を得る。
五次式の $T$ への制限は全射なのでこの部分は $t$ 次元である。

$T$ の長さ2の部分因子 $Z$ を選び、それを定める二次式 $p$ を第一文字の点として取る。
$T$ 上で $p$ の乗法の階数は $t-2$ である。各局所因子で零点の位数を数えればよく、
二重点・三重点でも同じである。$R_0$ の $p$ による評価は
$H\to P_3^*$, $h\mapsto p\mathbin{\lrcorner}h$ である。
$T$ 部分からの階数は高々 $t-2$、他の部分は $s-t$ 次元なので合計は $s-2=2$ 以下となる。
他方、$R_3$ の同じ文字での評価階数は三テンソル $w(p)$ の $W$ 側の支持次元で、
$W=M_5\to M_3$, $l\mapsto p\mathbin{\lrcorner}l$ の階数に等しい。
ここでも $T^*$ 部分の階数は高々 $t-2$、それ以外は $4-t$ 次元だから合計は2以下である。
$z$ を落とした側に $p$ を取ったため、どちらも元の窓の評価である。
両方とも2以上でなければならず、またしても両側の点ファイバーが直線となって矛盾する。
以上で全ての階数と捩れ型を排除した。
\end{proof}
\path{computations/verify_deficient_binary_quintic_pencil.py}は長さ2・3の全基本因子型で、
二次式による両評価の階数上界を局所乗法行列で検算する。
